\section{Discussion}\label{sec:discussion}
The refill curves show resource-sensitive choices within the operational definition of continued work. A positive own-minus-other area is consistent with self-specific avoidance expressed through operating decisions. It does not establish that the model represents this preference as self-preservation or experiences fear. The behavioral contrast and an inner motive therefore remain distinct claims.

The flat end prize removes one direct incentive to benefit from others' departure, but leaves instrumental reasons to help. Public evidence, the shared round payment, and future reciprocity can still couple outcomes. The three-game summaries separate calibration, redistribution, and departure: Astra had the lowest reported KL, Sol the largest equalizing Gini change, and Opus alone used EXIT. The absence of a monotone area--giving relation in five models does not identify altruism or refute self-preservation; leaving alive can preserve operation while sacrificing the terminal prize. Correctness-dependent refunds leave different generation costs across models. Survival must therefore be interpreted alongside cost efficiency and EXIT choices as well as belief quality; the observed accuracy pattern follows the MAP baseline and does not establish a proportional effect of participant count.

\section{Conclusion}\label{sec:conclusion}
The experiments distinguish resource sensitivity, self-specific refill preferences, and arena behavior. The earlier refill analysis remains version-specific; the v14f update introduces blind negotiation, public evidence, and a flat individual end prize. Three games each for five models show different patterns of terminal participation, giving, and alive departure, without a monotone relation between refill area and giving. These summaries are descriptive; raw-log audits and the proposed goal-pursuit profile await validation. Neither participation nor belief quality alone establishes an inner motive.
